% Capítulo 1 - Introducción

El presente informe corresponde al desarrollo del  proyecto de título para optar al grado de Ingeniero Civil en Informática de la Universidad del Bío-Bío. Este trabajo presenta DOMU, una plataforma integral para la administración y gestión de copropiedades que busca competir en el mercado chileno, ofreciendo una solución caracterizada por su alta integración, eficiencia operativa y accesibilidad económica.

En la actualidad, numerosas comunidades residenciales en Chile enfrentan desafíos críticos en su gestión diaria, tales como deficiencias en el control de acceso de visitas, falta de trazabilidad en la recepción de encomiendas, opacidad en la administración financiera y debilidades en la comunicación interna. Estas problemáticas no solo generan ineficiencia, sino que derivan en conflictos de convivencia y desconfianza por parte de los residentes. Este escenario se vuelve más complejo ante el crecimiento sostenido del parque inmobiliario y las exigencias de la nueva Ley de Copropiedad Inmobiliaria N\textsuperscript{o}~21.442, la cual establece estándares estrictos de transparencia, seguridad y profesionalismo en la administración.

Ante esta necesidad, surge DOMU, una plataforma que centraliza y automatiza los procesos clave del ecosistema residencial. A diferencia de las soluciones locales existentes que suelen abordar la gestión de forma parcial o fragmentada, DOMU integra las funciones administrativas, operativas y de seguridad en un solo sistema, involucrando activamente a todos los actores: administradores, comités, conserjes, residentes, proveedores y personal operativo.

Más allá de su componente técnico, DOMU tiene como propósito fortalecer la convivencia entre actores, facilitando el cumplimiento de la normativa vigente y promoviendo entornos residenciales más seguros, eficientes y participativos.


\newpage
\section{Presentación del área y su problemática}
Esta sección explica la situación actual e identifica los problemas que surgen de esta área.

\subsection{Description del problema}
La administración de comunidades residenciales en Chile enfrenta actualmente múltiples desafíos operativos, normativos y tecnológicos. El crecimiento sostenido del parque inmobiliario, junto con la entrada en vigor de la Ley de Copropiedad Inmobiliaria N\textsuperscript{o}~21.442 en abril de 2022, ha establecido nuevas exigencias en materia de trazabilidad, participación y transparencia en la gestión comunitaria. Sin embargo, en la práctica, muchas comunidades continúan funcionando con sistemas manuales y desarticulados, lo que afecta directamente la eficiencia del servicio y la calidad de vida de los residentes.

Entre las principales dificultades se encuentra el registro de visitas, que en la mayoría de los casos aún se realiza mediante libros físicos y anotaciones informales. Este método carece de mecanismos de validación y seguimiento, es susceptible a errores y representa un riesgo crítico frente a la Ley N\textsuperscript{o}~19.628 sobre Protección de la Vida Privada, al dejar información sensible (como nombres, RUT, entre otros) expuesta sin protocolos de cifrado, políticas de almacenamiento o control de acceso restringido.

Por otra parte la administración interna también presenta falencias importantes debido a la fragmentación de la data. La asignación de tareas al personal operativo (como personal de aseo, mayordomos, jardineros, etc) suele realizarse de forma informal (verbal), sin indicadores de desempeño (KPIs) ni seguimiento sistematizado. En el ámbito financiero, el uso de planillas de cálculo aisladas para la gestión de gastos comunes y morosidad dificulta la rendición de cuentas obligatoria que exige la ley de copropiedad, impidiendo una toma de decisiones basada en datos íntegros y actualizados.

En respuesta, han emergido diversas soluciones digitales en el mercado nacional que han liderado la digitalización del sector. No obstante, plataformas como \textit{ComunidadFeliz}, \textit{Edifito} y \textit{EdiPro} presentan modelos de servicio que imponen limitaciones críticas para una gestión integrada, operando mayoritariamente bajo un esquema de modularización comercial. En este modelo, funcionalidades esenciales (como la conciliación bancaria, el control de acceso, y otras) se encuentran segmentadas en diferentes niveles de suscripción, lo que obliga a las administraciones a combinar distintas herrmientas con el software para cumplir con las normas de la ley de copropiedad. Esta falta de integración no solo incrementa los costos operativos, sino que provoca una duplicidad de esfuerzos y una inconsistencia en la información almacenada, reduciendo la confiabilidad del sistema.

En este contexto, la falta de una solución digital unificada, económica y escalable evidencia un problema estructural en la gestión comunitaria. La dependencia de sistemas manuales y la fragmentación tecnológica dificultan la implementación de una administración transparente y eficiente. Esto justifica la necesidad de \textit{DOMU}, una propuesta que centraliza los procesos clave, garantiza la integridad de los datos y fortalece la convivencia en las comunidades residenciales bajo el marco normativo actual.

\section{Presentación de los procesos del ámbito del proyecto en la Empresa}

En esta sección se incluyen diagramas BPMN 2.0 de cuatro procesos que actualmente se realizan en la mayoria de comunidades que aun no se han digitalizado. Estos modelos servirán como punto de partida para el diseño y automatización de los flujos para nuestra plataforma \textit{DOMU}.

\subsection{Proceso de registro residente}

El proceso de registro comienza cuando un residente solicita ser inscrito mediante un formulario. La administración valida los datos ingresados. Si el formulario es inválido, se notifica el rechazo. Si es válido, se registra al residente y se le notifica. Finalmente, el residente valida su información, completando así su registro en una comunidad residencial.

\begin{figure}[H]
\centering
\includegraphics[width=0.95\linewidth]{Proyecto_de_Titulo/InformeFinal-SW/images/ProcesoDeRegistroResidenteBPMN.png}
\caption{Proceso de registro de residente}
\label{fig:bpmn-inicio-sesion}
\end{figure}

\subsection{Proceso de registro de visitas}

El proceso de registro de visitas, común en comunidades sin plataformas digitales o con planes básicos de estas, comienza con la llegada del visitante al condominio. El conserje notifica la visita al residente y solicita la confirmación para autorizar el ingreso. Paralelamente, se solicita al visitante que complete un formulario en papel con sus datos personales y motivo de la visita. Una vez entregado el formulario al conserje, este procede a validar la información registrada. Tras la validación y la confirmación del residente, el conserje permite el ingreso del visitante, finalizando el proceso con su entrada al condominio.

\begin{figure}[H]
\centering
\includegraphics[width=0.95\linewidth]{Proyecto_de_Titulo/InformeFinal-SW/images/ProcesoRegistroVisitaBPMN.jpg}
\caption{Proceso de registro de visitas en comunidades residenciales.}
\label{fig:bpmn-registro-visita}
\end{figure}

\subsection{Proceso de recepción y entrega de encomiendas}

El proceso de recepción de encomiendas, común en comunidades que no cuentan con plataformas digitales o utilizan planes básicos, comienza cuando un repartidor entrega el paquete a un conserje. Este registra manualmente la recepción, anotando los datos disponibles —como destinatario, remitente y hora de entrega— en un cuaderno o planilla física.

Luego de registrar la encomienda, el conserje notifica informalmente al residente, ya sea mediante una llamada telefónica, mensaje de texto o comunicación directa. El residente, al recibir la notificación, válida por su cuenta el envío con el proveedor para confirmar su origen y contenido. Una vez verificado, se dirige a retirar el paquete personalmente.

Finalmente, el residente firma un libro de recepción o documento similar como constancia de conformidad, completando así el proceso de entrega.


\begin{figure}[H]
\centering
\includegraphics[width=0.95\linewidth]{Proyecto_de_Titulo/InformeFinal-SW/images/ProcesoDeEncomiendasBPMN.jpg}
\caption{Proceso de encomiendas.}
\label{fig:bpmn-encomienda}
\end{figure}

\subsection{Proceso de pago de gastos comunes}

El proceso de pago de gastos comunes, tal como se realiza en muchas comunidades sin automatización digital, comienza con la llegada de la boleta mensual al residente, la cual suele ser enviada por correo electrónico o distribuida en formato físico (impresa). El residente revisa el detalle del cobro y, posteriormente, realiza una transferencia bancaria al número de cuenta indicado por la administración del edificio.

Una vez realizada la transferencia, el residente debe enviar manualmente el comprobante de pago —ya sea por correo electrónico, WhatsApp o entrega física— a la administración. Esta se encarga de validar el documento recibido, verificando el monto, la fecha y el nombre del remitente.

Si el pago es correcto, la administración actualiza el estado de los gastos comunes de manera manual en sus registros, y notifica al residente que su estado de cuenta ha sido actualizado. Finalmente, el residente puede comprobar que su boleta ha sido registrada como pagada, completando así el proceso.

\begin{figure}[H]
\centering
\includegraphics[width=0.95\linewidth]{Proyecto_de_Titulo/InformeFinal-SW/images/ProcedoDePagoGgCcBPMN.png}
\caption{Proceso de pago de gastos comunes.}
\label{fig:bpmn-pago-ggcc}
\end{figure}


\section{Propuesta de Solución}

Se propone una plataforma web que centralice los procesos críticos de la administración de copropiedades, desde el control de accesos hasta la gestión financiera y la comunicación interna.

A diferencia de las herramientas actuales, \textit{DOMU} integra un sistema de registro de visitas que permite a los residentes pre-autorizar visitantes, crear contactos frecuentes y agendar visitas con fecha y hora. La conserjería puede verificar el ingreso mediante códigos QR, validar el RUT y registrar el \textit{check-in}, facilitando el control y la trazabilidad de accesos sin depender de dispositivos físicos especializados.

Asimismo, la plataforma digitaliza el registro de encomiendas , la gestión de incidentes y reclamos con seguimiento por estado y asignación, y la reserva de áreas comunes con control de aforo, franjas horarias y bloqueo ante morosidad. El módulo financiero permite crear períodos de gastos comunes, cargar cargos, generar estados de cuenta en PDF y habilitar un flujo de pago que registra comprobantes y emite recibos.

La comunicación interna se apoya en un foro categorizado, un chat en tiempo real entre vecinos y una tienda comunitaria para compraventa dentro del edificio. Los residentes pueden participar en votaciones con cierre programado y exportación de resultados, y los administradores pueden gestionar personal, gestionar proveedores , asignar tareas y administrar unidades habitacionales y residentes desde un mismo entorno.

De este modo, la aplicación aborda de manera unificada la administración interna (tareas y personal), la experiencia del residente (pagos, reservas, comunicación, incidentes) y la transparencia financiera, promoviendo eficiencia y coordinación en la convivencia diaria.

\section{Análisis de los principales trabajos realizados en el área}

Actualmente existen varias soluciones orientadas a la administración de edificios y condominios, como \textit{ComunidadFeliz}, \textit{Edifito}, \textit{Edipro} y \textit{GastoComunChile}. Todas ellas abordan parcialmente la gestión financiera, la comunicación con residentes, el registro de gastos comunes y, en algunos casos, el control de accesos. Sin embargo, ninguna integra de forma completa las áreas involucradas —administrador, residentes, personal de aseo, conserjes, proveedores y visitas.

\begin{enumerate}
    \item \textbf{ComunidadFeliz}: ofrece una plataforma enfocada en la administración digital de comunidades residenciales, con funcionalidades que incluyen gestión de gastos comunes, pagos en línea, comunicación con residentes, votaciones virtuales, reservas de espacios comunes y conciliación bancaria. El Plan Full incorpora módulos adicionales como gestión de conserjería, remuneraciones y soporte especializado.

    \par\noindent
    Si bien \textit{ComunidadFeliz} cuenta con una presencia consolidada en el mercado chileno, su enfoque se centra principalmente en la administración financiera y la interacción con residentes, dejando fuera la integración operativa directa de actores como el personal interno, proveedores o sistemas físicos de acceso.

    \begin{figure}[H]
        \centering
        \includegraphics[width=.8\linewidth]{Proyecto_de_Titulo/InformeFinal-SW/images/PaginaComunidadFeliz.png}
        \caption{Página web de ComunidadFeliz.}
        \label{fig:site-comunidadfeliz}
    \end{figure}

    \item \textbf{Edifito}: ofrece una plataforma digital orientada a la administración integral de comunidades, con foco en la automatización de procesos financieros y operativos. Entre sus principales funcionalidades se encuentran la generación automática de gastos comunes, conciliación bancaria, gestión de remuneraciones, control de accesos mediante códigos QR o lectura de cédulas —similar a los lectores de precios de las cajas automáticas de supermercados—, administración de encomiendas, reservas de espacios comunes, comunicación interna, votaciones virtuales y conexión con cámaras de seguridad.

    \par\noindent
    \textit{Edifito} trabaja bajo un modelo comercial basado exclusivamente en cotización personalizada, por lo que no publica valores estándar en su sitio web. El precio final se ajusta en función del tamaño de la comunidad, la cantidad de unidades habitacionales y los módulos requeridos por cada cliente.

    \begin{figure}[H]
        \centering
        \includegraphics[width=.8\linewidth]{Proyecto_de_Titulo/InformeFinal-SW/images/PaginaEdifito.png}
        \caption{Página web de Edifito.}
        \label{fig:site-edifito}
    \end{figure}

    \item \textbf{Edipro}: ofrece una plataforma digital orientada a la administración de comunidades, centrada principalmente en la automatización de procesos financieros y operativos. Sus funcionalidades incluyen gestión de gastos comunes, conciliación bancaria, control de accesos mediante códigos QR o lectura de cédulas, administración de remuneraciones, registro de encomiendas, reservas de espacios comunes, votaciones virtuales, comunicación interna, videoconferencias e integración con cámaras de seguridad.

    \par\noindent
    \textit{Edipro} opera bajo un modelo modular, compuesto por tres paquetes escalables según las necesidades de cada comunidad. El \textit{Pack Conserje} (CLP 9.571 mensuales) contempla funciones básicas como control de accesos, registro de visitas y encomiendas, mensajería interna y tareas operativas. El \textit{Pack Administra} (CLP 25.524 mensuales) añade contabilidad, cálculo de gastos comunes, remuneraciones y conciliación bancaria. Finalmente, el \textit{Pack Full} (aproximadamente CLP 31.905 mensuales) integra todos los servicios anteriores, sumando reservas de espacios, videoconferencias y votaciones digitales.

    \par\noindent
    Estos valores se encuentran disponibles en el portal oficial:       \url{https://www.comparasoftware.cl/edipro}.

    \begin{figure}[H]
        \centering
        \includegraphics[width=.8\linewidth]{Proyecto_de_Titulo/InformeFinal-SW/images/PaginaEdipro.png}
        \caption{Página web de Edipro.}
        \label{fig:site-edipro}
    \end{figure}


    \item \textbf{GastoComunChile}: su propuesta se centra en optimizar la gestión financiera comunitaria, con funcionalidades que incluyen el ingreso de consumos y cobros individuales, cálculo automático de gastos comunes, conciliación bancaria, generación de minutas y reportes —como libro diario, mayor, cartolas de morosidad o corte de servicios—, envío masivo de notificaciones por correo y gestión de remuneraciones con archivos compatibles con \textit{Previred}.

    \par\noindent
    \textit{GastoComunChile} ofrece planes por comunidad con un costo inicial de \$25.000 CLP al mes (más IVA), o \$250.000 CLP al año si se prefiere pago anual, incluyendo soporte técnico, configuración personalizada y actualizaciones continuas sin cargo adicional.

    \par\noindent
    Aunque la plataforma cubre a fondo toda la gestión de gastos comunes y nómina, se posiciona como una competencia fuerte únicamente en ese ámbito específico.

    \begin{figure}[H]
        \centering
        \includegraphics[width=.8\linewidth]{Proyecto_de_Titulo/InformeFinal-SW/images/PaginaGastoComunChile.png}
        \caption{Página web de GastoComunChile.}
        \label{fig:site-gastocomunchile}
    \end{figure}

\end{enumerate}

\par\vspace{0.5em}
\noindent
\textit{DOMU} se diferencia al proponer un sistema integral que combina la plataforma de software con un sistema de acceso ágil, automatizando el registro de visitantes, proveedores y personal de mantenimiento. Además, incorpora la gestión centralizada de tareas, pagos, comunicación interna y votaciones en un solo entorno, solucionando la falta de trazabilidad y coordinación observada en los sistemas actuales.


% \section{Justificación de la Propuesta}

% La gestión de comunidades residenciales involucra una serie de procesos interdependientes, como el control de accesos, la administración financiera, la comunicación con residentes, la recepción de encomiendas y la coordinación del personal interno. Actualmente, muchas de estas tareas se llevan a cabo de forma fragmentada, ya sea mediante herramientas manuales o plataformas digitales aisladas, lo que genera ineficiencias, errores frecuentes, pérdida de tiempo y una limitada capacidad de control y trazabilidad.

% Las plataformas existentes han logrado digitalizar ciertos aspectos clave, sirviendo como referencia para identificar buenas prácticas y oportunidades de mejora. Sin embargo, ninguna de ellas integra todos los procesos críticos en un solo sistema ni contempla soluciones físicas que refuercen la seguridad operativa, como un tótem de registro. Esta ausencia de integración limita su efectividad a la hora de ofrecer una experiencia realmente centralizada, segura y eficiente para las comunidades.

% En este contexto, la propuesta de \textit{DOMU} se justifica por su enfoque integral e innovador. A través de una única plataforma, se reúnen todos los elementos esenciales de la administración comunitaria —pago de gastos comunes, remuneraciones, notificaciones, entre otros—. El uso de un tótem de acceso permite registrar y validar de manera inmediata tanto a visitantes como al personal, reduciendo errores humanos y fortaleciendo el control de ingreso. Además, el sistema permite asignar, monitorear y evaluar tareas internas, optimizando la productividad y la capacidad de respuesta operativa.

% \textit{DOMU} no solo digitaliza, sino que también centraliza y conecta todos los procesos clave de una comunidad, entregando una solución moderna, escalable y competitiva frente a las herramientas actuales. Su implementación apunta a mejorar la eficiencia administrativa, fortalecer la seguridad y promover una mejor convivencia en los entornos residenciales.

\section{Justificación de la Propuesta}

La gestión de comunidades residenciales involucra una serie de procesos interdependientes: el control de accesos, la administración financiera, la comunicación con residentes, la recepción de encomiendas y la coordinación del personal interno. En la práctica, muchas de estas tareas se realizan de forma fragmentada (mediante planillas, correos electrónicos o plataformas digitales aisladas), lo que genera ineficiencias, errores frecuentes, pérdida de tiempo y una limitada capacidad de control y trazabilidad. La desconexión entre módulos impide que administradores, conserjes y residentes compartan una misma visión del estado de la comunidad y de las acciones pendientes.

Las plataformas existentes han digitalizado aspectos específicos (gestión financiera, comunicación, reservas o accesos), sirviendo como referencia para identificar buenas prácticas. Sin embargo, la mayoría de ellas opera de forma modular o con enfoque parcial: algunas priorizan el área financiera y descuidan la operativa; otras facilitan el control de accesos pero no integran la gestión de encomiendas, incidentes o tareas internas. Esta ausencia de integración limita su efectividad para ofrecer una experiencia realmente centralizada, segura y eficiente en la comunidad.

En este contexto, la propuesta de \textit{DOMU} se justifica por su enfoque integral, centrado en una única plataforma que conecta los procesos críticos de la administración comunitaria. En lugar de depender de dispositivos físicos especializados, el sistema aprovecha el computador de conserjería y la web para permitir que los residentes pre-autoricen visitas, gestionen contactos frecuentes y que la portería verifique el ingreso mediante códigos QR y validación de RUT. Esta aproximación reduce errores de transcripción y mejora la trazabilidad sin requerir hardware adicional.

El control de accesos se complementa con la gestión de encomiendas (recepción, asignación a unidad y cambio de estado hasta la entrega), la gestión de incidentes y reclamos con seguimiento por estado y asignación, y la reserva de áreas comunes con control de aforo, franjas horarias y bloqueo ante morosidad. El módulo financiero permite crear períodos de gastos comunes, cargar cargos, generar estados de cuenta en PDF y registrar pagos con comprobantes y recibos. La comunicación interna se apoya en un foro categorizado, un chat en tiempo real entre vecinos y una tienda comunitaria para compraventa dentro del edificio. Además, el sistema permite asignar y monitorear tareas del personal, realizar votaciones con cierre programado y administrar unidades habitacionales y residentes desde un mismo entorno.

\textit{DOMU} no solo digitaliza procesos aislados, sino que los centraliza y conecta en un solo sistema, entregando una solución moderna, escalable y competitiva frente a las herramientas actuales. Su implementación apunta a mejorar la eficiencia administrativa, fortalecer el control y la transparencia, y promover una mejor convivencia en los entornos residenciales.

\section{Objetivos del proyecto}

En este punto se presentan los objetivos que guiarán el desarrollo, junto con las actividades correspondientes para su cumplimiento. Se expone tanto el objetivo general como los objetivos específicos.

\subsection{Objetivo General del proyecto}

Desarrollar una plataforma web integral para la administración de comunidades residenciales que integre y optimice los procesos críticos: registro y verificación de visitas con pre-autorización y control de acceso, trazabilidad de encomiendas, gestión de gastos comunes con generación de estados de cuenta y registro de pagos, asignación y seguimiento de tareas al personal, reporte y seguimiento de incidentes, reserva de áreas comunes, y comunicación interna mediante foro, chat y marketplace comunitario. El sistema debe ofrecer un entorno centralizado, trazabilidad de la información y acceso continuo para administradores, conserjes y residentes.


\subsection{Objetivos Específicos y Descripción de las actividades}

A continuación, se detallan los objetivos específicos del proyecto junto con  la descripción de las actividades asociadas a cada uno de estos. Estas actividades permitirán organizar de manera estructurada las tareas necesarias para alcanzar los objetivos planteados, asegurando un desarrollo sistemático y organizado del proyecto.
\newpage
\begin{longtable}{|p{0.4\textwidth}|p{0.55\textwidth}|}
\hline
\textbf{Objetivo} & \textbf{Descripción de las actividades} \\
\hline
\endfirsthead

\hline
\textbf{Objetivo} & \textbf{Descripción de las actividades} \\
\hline
\endhead

\hline
\endfoot

\hline
1. Comparar los sistemas actuales de administración de edificios, identificando sus fortalezas y limitaciones, así como las percepciones de los usuarios, para definir los requerimientos funcionales de nuestra propuesta. &
\begin{enumerate}
    \item Investigar y documentar al menos tres sistemas de administración de edificios existentes en el mercado chileno.
    \item Identificar las características y funcionalidades clave de cada sistema, considerando también la experiencia de uso reportada por los usuarios finales.
    \item Analizar las fortalezas y debilidades de cada sistema en relación con las necesidades detectadas de los distintos tipos de usuarios (administradores, residentes, personal).
    \item Recoger y sistematizar retroalimentación de usuarios reales mediante encuestas o entrevistas estructuradas (mínimo 5 personas).
    \item Definir los requerimientos funcionales del nuevo sistema basándose en el análisis comparativo y en las necesidades expresadas por los usuarios.
\end{enumerate}
\\ \hline

2. Definir tecnologías y metodologías disponibles que permitan desarrollar el sistema de registro con escaneo QR, página web y backend. &
\begin{enumerate}
    \item Definir tecnologías para el frontend y backend.
    \item Investigar y seleccionar el hardware de escaneo QR (pistola lectora), evaluando su compatibilidad  y la estructura de datos arrojada por la cédula de identidad chilena.
    \item Evaluar metodologías de desarrollo de software (ágil, cascada, etc.) adecuadas para el proyecto y escoger la que mejor se ajuste a la propuesta.
    \item Investigar opciones para la integración \textit{plug-and-play} del escáner QR con el computador de conserjería, y definir el método de procesamiento de la cadena de texto obtenida.
\end{enumerate}
\\ \hline

3. Diseñar la estructura de software que incorpora módulos de control de acceso, seguimiento de encomiendas, asignación de tareas al personal y comunicación con los residentes, entre otros. &
\begin{enumerate}
    \item Definir la arquitectura general del software y sus diferentes módulos.
    \item Diseñar el modelo de datos para cada módulo.
    \item Diseñar los diagramas de procesos en notación BPMN y los modelos entidad–relación (MER) correspondientes a los flujos.
    \item Diseñar la interfaz de usuario (UI) para cada módulo.
\end{enumerate}
\\ \hline

4. Implementar un prototipo funcional que integre los módulos diseñados en una plataforma inicial. &
\begin{enumerate}
    \item Configurar el entorno de desarrollo y las herramientas de control de versiones.
    \item Implementar las funcionalidades de backend.
    \item Desarrollar el frontend garantizando la comunicación con el backend.
    \item Realizar pruebas unitarias y ajustes.
    \item Asegurar la correcta sincronización de datos entre módulos.
\end{enumerate}
\\ \hline

5. Integrar el software de registro con el hardware de escaneo QR, verificando la correcta captura, \textit{parseo} y envío de los datos de la cédula de identidad a la plataforma central. &
\begin{enumerate}
    \item Adquirir el hardware (pistola escáner) y configurar los \textit{drivers} necesarios para la correcta \textit{emulación del teclado} en el computador de conserjería.
    \item Implementar la lógica de \textit{parsing} (análisis) en el software para decodificar la cadena de texto arrojada por el escáner QR de la cédula chilena.
    \item Garantizar la sincronización y el envío seguro de los datos personales capturados a la base de datos central.
    \item Validar la precisión de la captura y el \textit{parsing} de los campos de la cédula (RUT, Nombre, Apellido).
    \item Realizar pruebas de rendimiento y robustez del proceso de escaneo.
\end{enumerate}
\\ \hline

6. Realizar pruebas finales, validación y ajuste del sistema DOMU. &
\begin{enumerate}
    \item Diseñar un plan de pruebas integral.
    \item Recopilar retroalimentación de usuarios finales.
    \item Corregir y optimizar el sistema según resultados.
    \item Validar cumplimiento de requerimientos definidos.
    \item Documentar resultados y dejar el sistema operativo.
\end{enumerate}
\\ \hline

\end{longtable}

\section{Composición del Informe}
Aquí va como se organiza el informe, mencionar los capítulos que abrodan las etapas del proyecto.
describiendo brevemente c/u 